\subsection{ガウス型}\label{節：ガウス型}
  ここまでの漸化式では,差を取る,定数を引く,別の数列に置き換えるなどの操作から,既知の形へ帰着させてきた.
  ガウス記号を含む場合も,この方針は変わらない.
  ただし,変形の手掛かりとして,整数性,不等式,余り,偶奇,数え上げなどが必要になる.
  ガウス型では,整数問題で用いる考え方が前面に現れるのである.

  ガウス記号があるというだけで,一つの解法が決まるわけではない.
  本項では,\textbf{ガウス記号が式の中で何をしているか}に注目して,10種類の使い方を整理する.
  各解説の直後に代表例題を一つ置いた.
  同じ問題が複数の使い方に関係することもあるので,分類は排他的な区分ではなく,着眼点を選ぶための見取り図として使おう.

  \subsubsection{下準備と分類の見取り図}\label{節：ガウス型の下準備}
    ガウス記号$\lfloor x\rfloor$または$[x]$は,実数$x$を超えない最大の整数を表し,床関数と呼ばれる.
    整数$k$について,
    \[
      \lfloor x\rfloor=k\iff k\le x<k+1
    \]
    である.\,$x\ge0$では小数点以下を切り捨てた値に一致するが,
    負数では0に近づく切り捨てと異なる.例えば$\lfloor-1.5\rfloor=-2$である.
    天井関数$\lceil x\rceil$は$x$以上の最小の整数であり,
    \[
      \lceil x\rceil=k\iff k-1<x\le k
    \]
    で定義される.

    \BookFigFloorGraph

    式を読むときは,床関数が\textbf{項の値,添字,付加項のどれに作用するか}を最初に確かめよう.
    そのうえで,以下の使い方を考える.
    \begin{enumerate}[label=(\arabic*)]
      \item 整数性・不等式で整数部分を確定する.
      \item 同じ値を取る区間を群としてまとめる.
      \item 余りを取り出し,有限個の値と周期を調べる.
      \item 床関数の差で境界・条件・繰り上がりを検出する.
      \item 丸め誤差と係数の関係を使う.
      \item 整数への丸めで不動点と停止を調べる.
      \item 添字を二つに分割する.
      \item 商と余りで整数の桁を読む.
      \item 小数部分を残し,回転として読む.
      \item 条件を満たす整数や格子点の個数を数える.
    \end{enumerate}
    最初の6種類は主に値やその変化,次の2種類は添字,
    最後の2種類は小数部分や個数という意味に注目するものである.
    一般項を求めるだけでなく,周期や止まる値,何を数えているかを調べることにも意味がある.

    なお,整数$k$について$\lfloor x+k\rfloor=\lfloor x\rfloor+k$だが,
    実数同士では$\lfloor x+y\rfloor=\lfloor x\rfloor+\lfloor y\rfloor$とは限らない.
    数項を計算する実験は着眼点を探すために使い,予想は不等式や帰納法などで確かめよう.

    \subsubsection{整数性と不等式で整数部分を確定する}\label{節：ガウス型の整数部分の確定}
    \bookterm{floor}{ガウス記号}を見たら,まず中にある数が整数かどうかを考えよう.
    整数$k$と実数$x$に対して,
    \[
      \lfloor k+x\rfloor=k+\lfloor x\rfloor
    \]
    が成り立つ.例えば,整数$a_n$に対して$\lfloor a_n+\sqrt3\rfloor=a_n+1$である.
    ただし,これを使う前に,初項から整数が保たれることを確かめる必要がある.

    中身が整数との和でなくても,\,$k\le X<k+1$を示せば$\lfloor X\rfloor=k$と分かる.
    平方根なら,先に隣り合う平方数の間に挟むとよい.
    \textbf{\bookterm{floor-function}{床関数}の値を確定してから,既知の\bookterm{recurrence}{漸化式}として解く}のがこの使い方である.

    \noindent \bookblocklink{example-gauss-role-1}{problem}{answer}{【例題】}\par\nobreak
    \begin{enumerate}[label=(\arabic*),parsep=3pt,beginpenalty=10000,format=\bookitemlink{example-gauss-role-1}{problem}{answer}]
      \item \bookterm{sequence}{数列}を次のように定める.\bookterm{general}{一般項}を求めよ.
\BookMathLayout{            \[
              a_1=1,\qquad
              a_{n+1}=\left\lfloor\sqrt{a_n^2+4a_n+5}\right\rfloor
              \quad(n\ge1).
            \]}{            \[
              \begin{aligned}
              a_1&=1,\\
              a_{n+1}&=\left\lfloor\sqrt{a_n^2+4a_n+5}\right\rfloor
              \\              &\quad(n\ge1).
              \end{aligned}
            \]}
    \end{enumerate}

    \noindent \bookblocklink{example-gauss-role-1}{hint}{problem}{【方針】}\par\nobreak
    \begin{enumerate}[label=(\arabic*),parsep=3pt,beginpenalty=10000,format=\bookitemlink{example-gauss-role-1}{hint}{problem}]
      \item 整数$a_n$に対して,平方根の中身を$(a_n+2)^2$と$(a_n+3)^2$で挟みます.
            まず各項が非負整数であることを確かめましょう.
    \end{enumerate}

    \noindent \bookblocklink{example-gauss-role-1}{answer}{problem}{【解答】}\par\nobreak
    \begin{enumerate}[label=(\arabic*),parsep=3pt,beginpenalty=10000,format=\bookitemlink{example-gauss-role-1}{answer}{problem}]
      \item \[
              a_n=2n-1\quad(n\ge1).
            \]
    \end{enumerate}

    \noindent \bookblocklink{example-gauss-role-1}{solution}{problem}{【解法】}\par\nobreak
    \begin{enumerate}[label=(\arabic*),parsep=3pt,beginpenalty=10000,format=\bookitemlink{example-gauss-role-1}{solution}{problem}]
      \item 初項は非負整数であり,非負整数$a_n$から次の項も非負整数になる.
            よって帰納法により,すべての項が非負整数である.
            非負整数$x$について,
            \[
              (x+2)^2\le x^2+4x+5<(x+3)^2
            \]
            が成り立つ.左の差は1,右の差は$2x+4>0$である.
            平方根を取ると,
            \[
              x+2\le\sqrt{x^2+4x+5}<x+3
            \]
            なので,\bookterm{floor-function}{床関数}の値は$x+2$である.
            従って$a_{n+1}=a_n+2$となり,\sectionref{節：等差型}に帰着して$a_n=2n-1$を得る.
    \end{enumerate}

    \subsubsection{同じ値を取る区間を群としてまとめる}\label{節：ガウス型の群数列}
    \bookterm{floor-function}{床関数}は,中身が同じ二つの連続整数の間にある間,同じ値を取る.
    例えば,正整数$k,j$に対して,
    \[
      \lfloor\sqrt k\rfloor=j
      \iff j^2\le k\le(j+1)^2-1
    \]
    だから,値$j$が$2j+1$項続く.
    このようなまとまりを一つの群として数えるのが,群\bookterm{sequence}{数列}の考え方である.
    実際に$\lfloor\sqrt k\rfloor$は,
    \[
      \underbrace{1,1,1}_{3\text{個}},\,
      \underbrace{2,2,2,2,2}_{5\text{個}},\,
      \underbrace{3,3,3,3,3,3,3}_{7\text{個}},\,\ldots
    \]
    と並ぶ.

    \bookterm{difference}{階差}にこの形が現れたら,\sectionref{節：階差型}と同じように足す.
    \textbf{群の境界と,最後の群のどこまでを足すか}を区別しよう.
    似た形でも,和の上端が$n$か$n-1$かで境界が変わる.

    \noindent \bookblocklink{example-gauss-role-2}{problem}{answer}{【例題】}\par\nobreak
    \begin{enumerate}[label=(\arabic*),parsep=3pt,beginpenalty=10000,format=\bookitemlink{example-gauss-role-2}{problem}{answer}]
      \item \bookterm{sequence}{数列}を次のように定める.\bookterm{general}{一般項}を求めよ.
            \[
              a_1=2,\qquad
              a_{n+1}=a_n+\lfloor\sqrt{n+2}\rfloor
              \quad(n\ge1).
            \]
            必要ならば,\,$m=\lfloor\sqrt{n+1}\rfloor$を用いてよい.
    \end{enumerate}

    \noindent \bookblocklink{example-gauss-role-2}{hint}{problem}{【方針】}\par\nobreak
    \begin{enumerate}[label=(\arabic*),parsep=3pt,beginpenalty=10000,format=\bookitemlink{example-gauss-role-2}{hint}{problem}]
      \item \bookterm{difference}{階差}を足し,\,$\lfloor\sqrt k\rfloor$の和にします.
            \,$k=1,2$で値がともに1であることを使い,和の範囲をそろえましょう.
    \end{enumerate}

    \noindent \bookblocklink{example-gauss-role-2}{answer}{problem}{【解答】}\par\nobreak
    \begin{enumerate}[label=(\arabic*),parsep=3pt,beginpenalty=10000,format=\bookitemlink{example-gauss-role-2}{answer}{problem}]
      \item \,$m=\lfloor\sqrt{n+1}\rfloor$とおくと,
            \[
              a_n=m(n+2)-\frac{m(m+1)(2m+1)}6.
            \]
    \end{enumerate}

    \noindent \bookblocklink{example-gauss-role-2}{solution}{problem}{【解法】}\par\nobreak
    \begin{enumerate}[label=(\arabic*),parsep=3pt,beginpenalty=10000,format=\bookitemlink{example-gauss-role-2}{solution}{problem}]
      \item \bookterm{difference}{階差}を足すと,
            \[
              a_n=2+\sum_{k=3}^{n+1}\lfloor\sqrt k\rfloor
                  =\sum_{k=1}^{n+1}\lfloor\sqrt k\rfloor.
            \]
            \,$n=1$では最初の和を空の和として0とする.
            \,$N=n+1,\ m=\lfloor\sqrt N\rfloor$とおく.
            第$m-1$群までは全項,第$m$群は$k=m^2$から$N$までの$N-m^2+1$項を足すので,
            \begin{align*}
              a_n&=\sum_{j=1}^{m-1}j(2j+1)+m(N-m^2+1)\\
                 &=\frac{(m-1)m(4m+1)}6\\
                 &\quad+m(N-m^2+1)\\
                 &=m(N+1)-\frac{m(m+1)(2m+1)}6.
            \end{align*}
            \,$N=n+1$を戻せば\bookterm{general}{一般項}を得る.\,$m=1$でも最初の和は0となり,初項と一致する.

            \noindent \textbf{【別解】}\\
            \,$\lfloor\sqrt k\rfloor$は,\,$j^2\le k$を満たす正整数$j$の個数でもある.
            同じ組$(j,k)$を$j$ごとに数えると,
            \begin{align*}
              a_n&=\sum_{j=1}^{m}(N-j^2+1)\\
                 &=m(N+1)-\frac{m(m+1)(2m+1)}6.
            \end{align*}
            群ごとにまとめる方法と,何を数えているかを読み直す方法がつながっている.
    \end{enumerate}

    \subsubsection{余りを取り出し,値の範囲と周期を調べる}\label{節：ガウス型の余りと周期}
    整数$X$と正整数$m$に対して,
    \[
      X-m\left\lfloor\frac Xm\right\rfloor
    \]
    は$X$を$m$で割った余りであり,\,$0,1,\ldots,m-1$のいずれかである.
    これは\bookterm{floor-function}{床関数}の定義から,この式が0以上$m$未満の整数になるためである.
    大きな数を計算する式に見えても,\textbf{値が有限個に制限される}ことに注目しよう.

    \,$x\equiv y\pmod m$は,\,$x-y$が$m$の倍数であることを表す合同式である.
    次の値が$n$にも依存するなら,\,$a_n$だけでなく$n$の余りも一緒に追う.
    同じ組に戻り,その組から次の組が一意に決まれば,以降は繰り返す.
    有限個の状態を追う場合でも,初項から繰り返すとは限らず,途中から周期的になることもある.

    \noindent \bookblocklink{example-gauss-role-3}{problem}{answer}{【例題】}\par\nobreak
    \begin{enumerate}[label=(\arabic*),parsep=3pt,beginpenalty=10000,format=\bookitemlink{example-gauss-role-3}{problem}{answer}]
      \item \bookterm{sequence}{数列}を次のように定める.\bookterm{general}{一般項}と最小の正の周期を求めよ.
            \[
              a_1=1,\qquad
              a_{n+1}=X_n-4\left\lfloor\frac{X_n}{4}\right\rfloor,
            \]
            \[
              X_n=a_n^2+n^2+n\quad(n\ge1).
            \]
    \end{enumerate}

    \noindent \bookblocklink{example-gauss-role-3}{hint}{problem}{【方針】}\par\nobreak
    \begin{enumerate}[label=(\arabic*),parsep=3pt,beginpenalty=10000,format=\bookitemlink{example-gauss-role-3}{hint}{problem}]
      \item 次の項は$X_n$を4で割った余りです.
            \,$(a_n,n\bmod4)$の組を追うと,同じ状態への戻りを確かめられます.
    \end{enumerate}

    \noindent \bookblocklink{example-gauss-role-3}{answer}{problem}{【解答】}\par\nobreak
    \begin{enumerate}[label=(\arabic*),parsep=3pt,beginpenalty=10000,format=\bookitemlink{example-gauss-role-3}{answer}{problem}]
      \item \[
              a_n=
              \begin{cases}
                1 & (n\equiv0,1\pmod4),\\
                3 & (n\equiv2,3\pmod4).
              \end{cases}
            \]
            最小の正の周期は4である.
    \end{enumerate}

    \noindent \bookblocklink{example-gauss-role-3}{solution}{problem}{【解法】}\par\nobreak
    \begin{enumerate}[label=(\arabic*),parsep=3pt,beginpenalty=10000,format=\bookitemlink{example-gauss-role-3}{solution}{problem}]
      \item 初項は整数.整数$a_n$から$X_n$も次の項も整数になるので,各項は整数である.
            第2項以降は$0\le a_n<4$であり,
            \[
              a_{n+1}\equiv a_n^2+n(n+1)\pmod4
            \]
            となる.順に追うと,
            \[
              a_1,a_2,a_3,a_4,a_5=1,3,3,1,1.
            \]
            第1項と第5項では値が1,添字の余りが1で一致する.
            以降も同じ組から同じ次の組が決まるので,帰納法により$a_{n+4}=a_n$である.
            従って【解答】の\bookterm{general}{一般項}が得られる.
            第1項と第4項の値だけは同じでも,添字の余りが違うことに注意しよう.

            最小周期を確認する.1と2は$a_1\ne a_2,\ a_1\ne a_3$より周期ではなく,
            3も$a_2\ne a_5$より周期ではない.4は周期なので,最小の正の周期は4である.
    \end{enumerate}

    \subsubsection{床関数の差で境界・条件・繰り上がりを検出する}\label{節：ガウス型の境界と条件}
    正整数$n$と整数$m\ge2$に対して,
    \[
      \left\lfloor\frac nm\right\rfloor-
      \left\lfloor\frac{n-1}{m}\right\rfloor
      =\begin{cases}
         1 & (n\text{が}m\text{の倍数}),\\
         0 & (\text{それ以外})
       \end{cases}
    \]
    となる.商の整数部分が増える境界を検出しているのである.
    同様に,\,$\lfloor\sqrt n\rfloor-\lfloor\sqrt{n-1}\rfloor$は$n$が平方数のときだけ1になる.
    \textbf{式の形で条件を表し,条件に応じて変化する\bookterm{recurrence}{漸化式}を作れる}という使い方である.

    二つの数の和でも,整数の境界を越えたかを調べられる.
    小数部分を$\{x\}=x-\lfloor x\rfloor$と書けば,
    \begin{align*}
      &\lfloor x+y\rfloor-\lfloor x\rfloor-\lfloor y\rfloor\\
      &\qquad=\lfloor\{x\}+\{y\}\rfloor\in\{0,1\}.
    \end{align*}
    これは小数部分の和で繰り上がりが起きたときだけ1になる.
    \bookterm{floor-function}{床関数}は一般には和に分配できないが,その差には意味がある.

    \noindent \bookblocklink{example-gauss-role-4}{problem}{answer}{【例題】}\par\nobreak
    \begin{enumerate}[label=(\arabic*),parsep=3pt,beginpenalty=10000,format=\bookitemlink{example-gauss-role-4}{problem}{answer}]
      \item \bookterm{sequence}{数列}を次のように定める.\bookterm{general}{一般項}を求め,\bookterm{difference}{階差}の意味を説明せよ.
            \[
              a_0=0,\qquad a_n=a_{n-1}+d_n\quad(n\ge1),
            \]
            \[
              \begin{aligned}
                d_n={}&\left\lfloor\frac n4\right\rfloor
                       -\left\lfloor\frac{n-1}{4}\right\rfloor\\
                      &-\left\lfloor\frac n6\right\rfloor
                       +\left\lfloor\frac{n-1}{6}\right\rfloor.
              \end{aligned}
            \]
    \end{enumerate}

    \noindent \bookblocklink{example-gauss-role-4}{hint}{problem}{【方針】}\par\nobreak
    \begin{enumerate}[label=(\arabic*),parsep=3pt,beginpenalty=10000,format=\bookitemlink{example-gauss-role-4}{hint}{problem}]
      \item 各差を倍数かどうかの判定として読みます.
            和を取ると,途中の\bookterm{floor-function}{床関数}が打ち消し合います.
    \end{enumerate}

    \noindent \bookblocklink{example-gauss-role-4}{answer}{problem}{【解答】}\par\nobreak
    \begin{enumerate}[label=(\arabic*),parsep=3pt,beginpenalty=10000,format=\bookitemlink{example-gauss-role-4}{answer}{problem}]
      \item \[
              a_n=\left\lfloor\frac n4\right\rfloor
                  -\left\lfloor\frac n6\right\rfloor.
            \]
            \,$d_n$は4の倍数だけなら1,6の倍数だけなら$-1$,両方またはいずれでもないなら0である.
    \end{enumerate}

    \noindent \bookblocklink{example-gauss-role-4}{solution}{problem}{【解法】}\par\nobreak
    \begin{enumerate}[label=(\arabic*),parsep=3pt,beginpenalty=10000,format=\bookitemlink{example-gauss-role-4}{solution}{problem}]
      \item \bookterm{floor-function}{床関数}の差の性質から,\,$d_n$は【解答】に述べた条件を表す.
            また,\,$a_n=\sum_{k=1}^{n}d_k$であり,
            \begin{align*}
              \sum_{k=1}^{n}\left(
                \left\lfloor\frac k4\right\rfloor-
                \left\lfloor\frac{k-1}{4}\right\rfloor\right)
                &=\left\lfloor\frac n4\right\rfloor,\\
              \sum_{k=1}^{n}\left(
                \left\lfloor\frac k6\right\rfloor-
                \left\lfloor\frac{k-1}{6}\right\rfloor\right)
                &=\left\lfloor\frac n6\right\rfloor.
            \end{align*}
            両式の差から\bookterm{general}{一般項}を得る.\,$n=0$でも両辺は0である.
            条件を一つずつ場合分けして値を追う方法と,差を足して一度に求める方法がある.
    \end{enumerate}

    \subsubsection{丸め誤差と係数の関係からガウス記号を外す}\label{節：ガウス型の丸め誤差}
    \,$a_{n+1}=\lfloor\lambda a_n\rfloor$では,整数との和の場合と違い,\,$\lambda$の整数部分を取り出すだけでは解けない.
    代わりに,
    \[
      \varepsilon_n=\lambda a_n-a_{n+1},
      \qquad 0\le\varepsilon_n<1
    \]
    とおき,丸めで捨てた部分を残して計算しよう.
    \textbf{誤差そのものではなく,誤差の範囲から次の整数部分を確定する}のが狙いである.

    係数$\lambda$が2次方程式を満たす場合は,\,$\lambda^2$を$\lambda$と整数で表せることが手掛かりになる.
    ただし,最後に現れる誤差がどの整数の間にあるかを示す必要がある.
    どの無理数の係数でも同じ隣接3項間漸化式へ帰着できるわけではない.

    \noindent \bookblocklink{example-gauss-role-5}{problem}{answer}{【例題】}
    \begin{enumerate}[label=(\arabic*),parsep=3pt,beginpenalty=10000,format=\bookitemlink{example-gauss-role-5}{problem}{answer}]
      \item 数列を次のように定める.ガウス記号を含まない漸化式に帰着させ,一般項を求めよ.
            \[
              a_1=1,\qquad
              a_{n+1}=\lfloor(1+\sqrt2)a_n\rfloor
              \quad(n\ge1).
            \]
    \end{enumerate}

    \noindent \bookblocklink{example-gauss-role-5}{hint}{problem}{【方針】}
    \begin{enumerate}[label=(\arabic*),parsep=3pt,beginpenalty=10000,format=\bookitemlink{example-gauss-role-5}{hint}{problem}]
      \item \,$\lambda=1+\sqrt2$とおき,\,$\lambda^2=2\lambda+1$を使います.
            整数から正の1未満の誤差を引く形にして,床関数の値を確定しましょう.
    \end{enumerate}

    \noindent \bookblocklink{example-gauss-role-5}{answer}{problem}{【解答】}
    \begin{enumerate}[label=(\arabic*),parsep=3pt,beginpenalty=10000,format=\bookitemlink{example-gauss-role-5}{answer}{problem}]
      \item \[
              a_1=1,\quad a_2=2,\qquad
              a_{n+2}=2a_{n+1}+a_n-1.
            \]
            \,$\alpha=1+\sqrt2,\ \beta=1-\sqrt2$とおくと,
            \[
              a_n=\frac12+\frac{\alpha^n+\beta^n}{4}.
            \]
    \end{enumerate}

    \noindent \bookblocklink{example-gauss-role-5}{solution}{problem}{【解法】}
    \begin{enumerate}[label=(\arabic*),parsep=3pt,beginpenalty=10000,format=\bookitemlink{example-gauss-role-5}{solution}{problem}]
      \item \,$\lambda=1+\sqrt2>2$とおく.正整数$a_n$から次の項も正整数となる.
            また,\,$\lambda a_n$は無理数なので,
            \[
              \varepsilon_n=\lambda a_n-a_{n+1}
            \]
            は$0<\varepsilon_n<1$を満たす.
            \,$\lambda^2=2\lambda+1$を使うと,
            \begin{align*}
              \lambda a_{n+1}
                &=\lambda(\lambda a_n-\varepsilon_n)\\
                &=2a_{n+1}+a_n-(\lambda-2)\varepsilon_n.
            \end{align*}
            \,$0<\lambda-2=\sqrt2-1<1$なので,最後に引く数も0と1の間にある.
            \,$M_n=2a_{n+1}+a_n$は整数だから,
            \[
              M_n-1<\lambda a_{n+1}<M_n.
            \]
            従って$a_{n+2}=M_n-1$となり,【解答】の隣接3項間漸化式を得る.
            初期条件は$a_2=\lfloor1+\sqrt2\rfloor=2$である.

            定数項を消すため$b_n=a_n-\frac12$とおくと,
            \[
              b_{n+2}=2b_{n+1}+b_n,
              \qquad b_1=\frac12,\quad b_2=\frac32.
            \]
            \sectionref{節：隣接3項間漸化式型}で学んだ特性方程式は$x^2-2x-1=0$で,
            その解が$\alpha,\beta$である.
            \,$\alpha+\beta=2,\ \alpha^2+\beta^2=6$より,
            \[
              b_n=\frac{\alpha^n+\beta^n}{4}
            \]
            は初期条件と漸化式を満たす.
            二つの初期値から各項は一意に決まるので,この式が一般項である.
            丸め誤差を評価した後は,既習の線形の漸化式に戻ることができた.
    \end{enumerate}


    \subsubsection{整数への丸めで不動点と停止を調べる}\label{節：ガウス型の不動点と停止}
    丸めを繰り返す数列では,一般項だけでなく,どの整数で止まるかにも注目しよう.
    次の値が今の値と等しくなる値を不動点という.
    不動点に到達すると,その後は同じ値が続く.
    ただし,丸めを外した実数の漸化式と,不動点が同じとは限らない.

    整数$K$を引いて,
    \[
      b_{n+1}=\left\lfloor\frac{b_n}{2}\right\rfloor
    \]
    に帰着できる場合がある.
    正整数$m$に対して,
    \[
      \left\lfloor\frac{\lfloor x\rfloor}{m}\right\rfloor
      =\left\lfloor\frac xm\right\rfloor
    \]
    が成り立つので,繰り返すと分母をまとめられる.
    この性質を確かめよう.\,$k=\lfloor x\rfloor$を$m$で割り,\,$k=mq+r$と書くと,
    \,$0\le r\le m-1$であり,\,$x=mq+r+t\ (0\le t<1)$となる.
    これより$0\le r+t<m$だから,どちらの床関数も$q$になる.

    \textbf{初項の位置と,負の数の床関数の値}も重要である.
    例えば,\,$-1<x<0$なら$\lfloor x\rfloor=-1$であって0ではない.

    \noindent \bookblocklink{example-gauss-role-6}{problem}{answer}{【例題】}
    \begin{enumerate}[label=(\arabic*),parsep=3pt,beginpenalty=10000,format=\bookitemlink{example-gauss-role-6}{problem}{answer}]
      \item 非負整数$M$に対し,数列を次のように定める.
            \[
              a_1=M,\qquad
              a_{n+1}=\left\lfloor\frac{a_n+6}{2}\right\rfloor
              \quad(n\ge1).
            \]
            一般項と,最終的に一定になる値を求めよ.
    \end{enumerate}

    \noindent \bookblocklink{example-gauss-role-6}{hint}{problem}{【方針】}
    \begin{enumerate}[label=(\arabic*),parsep=3pt,beginpenalty=10000,format=\bookitemlink{example-gauss-role-6}{hint}{problem}]
      \item 6を引くと,床関数によって半分にする漸化式になります.
            初項が6より小さい場合は,負の数の床関数を使って考えましょう.
    \end{enumerate}

    \noindent \bookblocklink{example-gauss-role-6}{answer}{problem}{【解答】}
    \begin{enumerate}[label=(\arabic*),parsep=3pt,beginpenalty=10000,format=\bookitemlink{example-gauss-role-6}{answer}{problem}]
      \item \[
              a_n=6+\left\lfloor\frac{M-6}{2^{n-1}}\right\rfloor.
            \]
            有限回の操作の後で,\,$M<6$なら5,\,$M\ge6$なら6に固定される.
    \end{enumerate}

    \noindent \bookblocklink{example-gauss-role-6}{solution}{problem}{【解法】}
    \begin{enumerate}[label=(\arabic*),parsep=3pt,beginpenalty=10000,format=\bookitemlink{example-gauss-role-6}{solution}{problem}]
      \item 初項からすべての項が非負整数となる.
            \,$b_n=a_n-6$とおくと,整数6を床関数の外に出せるので,
            \begin{align*}
              b_{n+1}
                &=\left\lfloor\frac{b_n+12}{2}\right\rfloor-6\\
                &=\left\lfloor\frac{b_n}{2}\right\rfloor.
            \end{align*}
            床関数の合成の性質と帰納法により,
            \[
              b_n=\left\lfloor\frac{M-6}{2^{n-1}}\right\rfloor.
            \]
            \,$M\ge6$では,十分大きな$n$で中身が$[0,1)$に入り,床は0になる.
            \,$M<6$では,十分大きな$n$で中身が$(-1,0)$に入り,床は$-1$になる.
            これより【解答】の値に有限回で到達する.

            実際,整数$x$が不動点となる条件は,
            \[
              x\le\frac{x+6}{2}<x+1
              \iff 4<x\le6
            \]
            なので,不動点は5と6の二つである.
            初項によって到達する不動点が分かれることは,この後の\sectionref{節：初項依存型}ともつながる.
    \end{enumerate}


    \subsubsection{床関数で添字を二つに分割する}\label{節：ガウス型の添字の分割}
    ここからは,\bookterm{floor-function}{床関数}が\bookterm{sequence}{数列}の値ではなく,\textbf{どの項を参照するか}を指定する使い方を考える.
    正整数$n$に対して,
    \[
      \left\lfloor\frac{n+1}{2}\right\rfloor
      +\left\lfloor\frac n2\right\rfloor=n
    \]
    であり,二つの添字は$n$をできるだけ等しく分けた大きさになる.
    偶数では両方が同じ添字になり,奇数では隣り合う添字になることが手掛かりである.

    この形の問題では,まず偶奇で分けて\bookterm{floor-function}{床関数}を外し,それから\bookterm{difference}{階差}を取る方法がある.
    半分の添字との関係が得られれば,\,$1,\ 2\le n<4,\ 4\le n<8,\ldots$のように,
    2の累乗ごとの群で整理できる.
    定義の確認も必要で,例えば$n=1$で元と同じ添字を参照する式は,
    そのまま初項を決める式として使えないことがある.

    \noindent \bookblocklink{example-gauss-role-7}{problem}{answer}{【例題】}\par\nobreak
    \begin{enumerate}[label=(\arabic*),parsep=3pt,beginpenalty=10000,format=\bookitemlink{example-gauss-role-7}{problem}{answer}]
      \item \bookterm{sequence}{数列}を次のように定める.\bookterm{general}{一般項}を求めよ.
\BookMathLayout{            \[
              a_1=2,\qquad
              a_n=a_{\left\lfloor\frac{n+1}{2}\right\rfloor}
                  +a_{\left\lfloor\frac n2\right\rfloor}+n-1
              \quad(n\ge2).
            \]}{            \[
              \begin{aligned}
              a_1&=2,\\
              a_n&=a_{\left\lfloor\frac{n+1}{2}\right\rfloor}
                  +a_{\left\lfloor\frac n2\right\rfloor}+n-1
              \\              &\quad(n\ge2).
              \end{aligned}
            \]}
            必要ならば,\,$2^m\le n<2^{m+1}$を満たす整数$m\ge0$を用いてよい.
    \end{enumerate}

    \noindent \bookblocklink{example-gauss-role-7}{hint}{problem}{【方針】}\par\nobreak
    \begin{enumerate}[label=(\arabic*),parsep=3pt,beginpenalty=10000,format=\bookitemlink{example-gauss-role-7}{hint}{problem}]
      \item 偶奇で式を分け,\,$b_n=a_{n+1}-a_n$を考えます.
            \,$b_{2k}$と$b_{2k+1}$が$b_k+1$になることを示し,群ごとに足しましょう.
    \end{enumerate}

    \noindent \bookblocklink{example-gauss-role-7}{answer}{problem}{【解答】}\par\nobreak
    \begin{enumerate}[label=(\arabic*),parsep=3pt,beginpenalty=10000,format=\bookitemlink{example-gauss-role-7}{answer}{problem}]
      \item \,$2^m\le n<2^{m+1}$を満たす整数$m\ge0$に対して,
            \[
              a_n=(m+3)n-2^{m+1}+1.
            \]
    \end{enumerate}

    \noindent \bookblocklink{example-gauss-role-7}{solution}{problem}{【解法】}\par\nobreak
    \begin{enumerate}[label=(\arabic*),parsep=3pt,beginpenalty=10000,format=\bookitemlink{example-gauss-role-7}{solution}{problem}]
      \item \,$n\ge2$では二つの添字は1以上$n-1$以下なので,\bookterm{sequence}{数列}は順に定義できる.
            偶奇で分けると,\,$k\ge1$について,
            \[
              \begin{aligned}
                a_{2k}&=2a_k+2k-1,\\
                a_{2k+1}&=a_{k+1}+a_k+2k.
              \end{aligned}
            \]
            \,$b_n=a_{n+1}-a_n$とおくと,両式の差から$b_{2k}=b_k+1$.
            また,\,$a_{2k+2}=2a_{k+1}+2k+1$を使うと$b_{2k+1}=b_k+1$である.
            \,$a_2=5$より$b_1=3$なので,
            \[
              b_{2k}=b_{2k+1}=b_k+1,\qquad b_1=3.
            \]

            \,$2^m\le n<2^{m+1}$の範囲では$b_n=m+3$となる.
            \,$m=0$では$b_1=3$と一致する.
            次の群の添字を半分にして床を取ると前の群に入り,値が1増えるから,
            群番号$m$についての帰納法でこの式が示される.

            従って,完全な群と最後の群を分けて足すと,
            \begin{align*}
              a_n={}&2+\sum_{j=0}^{m-1}2^j(j+3)\\
                   &+(n-2^m)(m+3).
            \end{align*}
            ここで,
            \[
              \sum_{j=0}^{m-1}2^j=2^m-1,
            \]
            \[
              \sum_{j=0}^{m-1}j2^j=(m-2)2^m+2
            \]
            を使う.後者は,和を$T$とおき$2T-T$を計算すれば得られる.
            \,$m=0$でも空の和として両式は成り立つ.
            これを代入すると,
            \[
              a_n=(m+3)n-2^{m+1}+1
            \]
            を得る.添字を分割する問題から,\bookterm{difference}{階差}を通じて群\bookterm{sequence}{数列}へ戻ることができた.
    \end{enumerate}

    \subsubsection{商と余りで整数の桁を読む}\label{節：ガウス型の桁の抽出}
    正整数$b\ge2$に対し,非負整数$n$を$b$で割って,
    \[
      q=\left\lfloor\frac nb\right\rfloor,
      \qquad r=n-b\left\lfloor\frac nb\right\rfloor
    \]
    とおくと,\,$n=bq+r,\ 0\le r<b$である.
    十進法では,\,$q=\lfloor\frac n{10}\rfloor$が一の位を取り除いた整数,
    \,$r=n-10\lfloor\frac n{10}\rfloor$が一の位の数字になる.

    一般の$b$進法も,各桁の数字を0から$b-1$までとし,位を$b$倍ずつにする表し方である.
    このとき$q$は最後の桁を取り除く操作,$r$はその桁の数字に当たる.
    \textbf{添字を小さくする床関数と,桁を取り出す余りを組み合わせる}ことができる.
    前の添字分割は二つの項を足したが,ここでは一つの項へ戻りながら最後の桁を処理する.

    \noindent \bookblocklink{example-gauss-role-8}{problem}{answer}{【例題】}
    \begin{enumerate}[label=(\arabic*),parsep=3pt,beginpenalty=10000,format=\bookitemlink{example-gauss-role-8}{problem}{answer}]
      \item 数列を次のように定める.
            \[
              a_0=0,\qquad
              a_n=a_{\left\lfloor\frac n3\right\rfloor}
                  +n-3\left\lfloor\frac n3\right\rfloor
              \quad(n\ge1).
            \]
            \,$a_n$が$n$の3進表示とどのような関係を持つか説明し,\,$a_{14}$を求めよ.
    \end{enumerate}

    \noindent \bookblocklink{example-gauss-role-8}{hint}{problem}{【方針】}
    \begin{enumerate}[label=(\arabic*),parsep=3pt,beginpenalty=10000,format=\bookitemlink{example-gauss-role-8}{hint}{problem}]
      \item \,$n=3q+r\ (0\le r<3)$と書くと,漸化式は$a_n=a_q+r$になります.
            \,$q$の3進表示に最後の桁$r$を付けると$n$になることを使いましょう.
    \end{enumerate}

    \noindent \bookblocklink{example-gauss-role-8}{answer}{problem}{【解答】}
    \begin{enumerate}[label=(\arabic*),parsep=3pt,beginpenalty=10000,format=\bookitemlink{example-gauss-role-8}{answer}{problem}]
      \item \,$a_n$は$n$の3進表示の各桁の数字の和である.
            \,$14=(112)_3$より,
            \[
              a_{14}=1+1+2=4.
            \]
    \end{enumerate}

    \noindent \bookblocklink{example-gauss-role-8}{solution}{problem}{【解法】}
    \begin{enumerate}[label=(\arabic*),parsep=3pt,beginpenalty=10000,format=\bookitemlink{example-gauss-role-8}{solution}{problem}]
      \item \,$n\ge1$では$0\le q=\lfloor\frac n3\rfloor<n$なので,初期値$a_0$から各項を定義できる.
            \,$r=n-3q$とおけば$r$は0,1,2のいずれかで,
            \[
              a_n=a_q+r.
            \]
            \,$n$より小さな整数について桁の和との一致が分かっているとしよう.
            \,$n=3q+r$の3進表示は,\,$q$の表示を一桁ずらして最後に$r$を付けたものになる.
            その桁の和は$q$の桁の和に$r$を足した値であり,漸化式と一致する.
            \,$a_0=0$から,この考えを$n=1,2,\ldots$と順に適用すれば,すべての$n$について示される.
            \,$q=0$の場合も,桁の和を0として同じ議論が成り立つ.

            実際,\,$14=3\cdot4+2,\ 4=3\cdot1+1,\ 1=3\cdot0+1$なので,
            \[
              a_{14}=a_4+2=a_1+3=a_0+4=4.
            \]
            この問題では,一般項を長い式で書く前に,数列が何を表すかを理解することに意味がある.
    \end{enumerate}


    \subsubsection{小数部分を残し,区間内の回転として読む}\label{節：ガウス型の小数部分と回転}
    \,$\{x\}=x-\lfloor x\rfloor$は$x$の小数部分を表し,\,$0\le\{x\}<1$である.
    例えば,\,$\{-0.2\}=0.8$となる.
    ある量を加えてから整数部分を捨てる操作は,値を$[0,1)$へ戻す操作である.
    区間の両端をつないで円周と見れば,毎回同じ割合だけ進む回転として解釈できる.

    \textbf{小数部分を追うと,整数の境界を越えても運動の規則が残る}.
    途中で何回境界を越えたかは,\bookterm{floor-function}{床関数}で記録できる.
    有理数の割合なら繰り返しが生じるが,無理数の場合については,周期があるかを別に証明しよう.

    \noindent \bookblocklink{example-gauss-role-9}{problem}{answer}{【例題】}\par\nobreak
    \begin{enumerate}[label=(\arabic*),parsep=3pt,beginpenalty=10000,format=\bookitemlink{example-gauss-role-9}{problem}{answer}]
      \item \,$\alpha=\sqrt2-1$とし,\bookterm{sequence}{数列}を次のように定める.
            \[
              a_0=0,\qquad
              a_{n+1}=a_n+\alpha-\lfloor a_n+\alpha\rfloor
              \quad(n\ge0).
            \]
            \bookterm{general}{一般項}を求め,異なる添字の項が同じ値になるか調べよ.
    \end{enumerate}

    \noindent \bookblocklink{example-gauss-role-9}{hint}{problem}{【方針】}\par\nobreak
    \begin{enumerate}[label=(\arabic*),parsep=3pt,beginpenalty=10000,format=\bookitemlink{example-gauss-role-9}{hint}{problem}]
      \item 整数部分を捨てると小数部分だけが残ることから,\,$n\alpha$の小数部分を予想します.
            一致する二項があれば,添字の差と$\alpha$の積が整数になることを使いましょう.
    \end{enumerate}

    \noindent \bookblocklink{example-gauss-role-9}{answer}{problem}{【解答】}\par\nobreak
    \begin{enumerate}[label=(\arabic*),parsep=3pt,beginpenalty=10000,format=\bookitemlink{example-gauss-role-9}{answer}{problem}]
      \item \[
              a_n=n\alpha-\lfloor n\alpha\rfloor\quad(n\ge0).
            \]
            異なる添字の項は同じ値にならず,周期的ではない.
    \end{enumerate}

    \noindent \bookblocklink{example-gauss-role-9}{solution}{problem}{【解法】}\par\nobreak
    \begin{enumerate}[label=(\arabic*),parsep=3pt,beginpenalty=10000,format=\bookitemlink{example-gauss-role-9}{solution}{problem}]
      \item \,$n=0$では予想した式は$a_0=0$を与える.
            \,$a_n=n\alpha-\lfloor n\alpha\rfloor$とすると,整数との和の性質から,
            \begin{align*}
              \lfloor a_n+\alpha\rfloor
                &=\lfloor(n+1)\alpha\rfloor-\lfloor n\alpha\rfloor,\\
              a_{n+1}
                &=(n+1)\alpha-\lfloor(n+1)\alpha\rfloor.
            \end{align*}
            よって帰納法により\bookterm{general}{一般項}が示された.すべての項は$[0,1)$に属する.

            \,$a_{n+p}=a_n\ (p\ge1)$とすると,
            \[
              p\alpha=\lfloor(n+p)\alpha\rfloor-\lfloor n\alpha\rfloor
            \]
            が整数になり,\,$\alpha$が有理数となる.これは$\alpha=\sqrt2-1$が無理数であることに矛盾する.
            従って異なる添字の項は一致しない.

            同じ\bookterm{recurrence}{漸化式}で$\alpha$が有理数$\frac pq\ (0<p<q)$であり,\,$p,q$が互いに素なら,
            最小の正の周期は$q$となる.\,$q\alpha$が整数なので$q$項後に戻り,
            戻るまでの項数$T$は$T\alpha$が整数となる条件から$q$の倍数だからである.
            回転の見方と,有理数・無理数という整数問題の考え方が結びついている.
    \end{enumerate}

    \subsubsection{床関数を整数の個数として数える}\label{節：ガウス型の個数と格子点}
    \,$X\ge0$に対して$\lfloor X\rfloor$は,\,$1\le k\le X$を満たす整数$k$の個数である.
    例えば$\lfloor\frac nm\rfloor$は,\,$1,\ldots,n$に含まれる$m$の倍数の個数を表す.
    \bookterm{floor-function}{床関数}を含む和は,\textbf{何を数えているか}から意味を読み直せる.

    座標がともに整数である点を格子点という.
    例えば$\sum_{k=1}^{n}\lfloor ck\rfloor\ (c>0)$は,
    \,$x=k$を固定するたびに,\,$1\le y\le ck$を満たす整数$y$を数えている.
    同じ値の群で整理する方法とは別に,図形の中の点を数える方法としても理解できる.
    二つの使い方は重なり,前の群\bookterm{sequence}{数列}の別解も整数の組を数え直す方法であった.

    \noindent \bookblocklink{example-gauss-role-10}{problem}{answer}{【例題】}\par\nobreak
    \begin{enumerate}[label=(\arabic*),parsep=3pt,beginpenalty=10000,format=\bookitemlink{example-gauss-role-10}{problem}{answer}]
      \item \bookterm{sequence}{数列}を次のように定める.
            \[
              a_0=0,\qquad
              a_n=a_{n-1}+\left\lfloor\frac{5n}{2}\right\rfloor
              \quad(n\ge1).
            \]
            \bookterm{general}{一般項}を求め,\,$a_n$をある範囲の格子点の個数として説明せよ.
    \end{enumerate}

    \noindent \bookblocklink{example-gauss-role-10}{hint}{problem}{【方針】}\par\nobreak
    \begin{enumerate}[label=(\arabic*),parsep=3pt,beginpenalty=10000,format=\bookitemlink{example-gauss-role-10}{hint}{problem}]
      \item \,$\lfloor\frac{5k}{2}\rfloor=2k+\lfloor\frac k2\rfloor$と分けます.
            後者の和を偶奇で求め,\bookterm{floor-function}{床関数}が数えている整数$y$の範囲も確認しましょう.
    \end{enumerate}

    \noindent \bookblocklink{example-gauss-role-10}{answer}{problem}{【解答】}\par\nobreak
    \begin{enumerate}[label=(\arabic*),parsep=3pt,beginpenalty=10000,format=\bookitemlink{example-gauss-role-10}{answer}{problem}]
      \item \[
              a_n=n(n+1)+\left\lfloor\frac{n^2}{4}\right\rfloor.
            \]
            これは整数$x,y$で,
            \[
              1\le x\le n,\qquad 1\le y\le\frac{5x}{2}
            \]
            を満たす格子点の個数である.
    \end{enumerate}

    \noindent \bookblocklink{example-gauss-role-10}{solution}{problem}{【解法】}\par\nobreak
    \begin{enumerate}[label=(\arabic*),parsep=3pt,beginpenalty=10000,format=\bookitemlink{example-gauss-role-10}{solution}{problem}]
      \item \bookterm{difference}{階差}を足し,整数$2k$を\bookterm{floor-function}{床関数}の外に出すと,
            \begin{align*}
              a_n&=\sum_{k=1}^{n}\left\lfloor\frac{5k}{2}\right\rfloor\\
                 &=n(n+1)+\sum_{k=1}^{n}\left\lfloor\frac k2\right\rfloor.
            \end{align*}
            後ろの和は$n=2m$なら,
            \[
              0+1+1+2+2+\cdots+(m-1)+(m-1)+m=m^2,
            \]
            \,$n=2m+1$なら$m^2+m$である.
            従っていずれの場合も,
            \[
              \sum_{k=1}^{n}\left\lfloor\frac k2\right\rfloor
              =\left\lfloor\frac{n^2}{4}\right\rfloor
            \]
            となり,\bookterm{general}{一般項}を得る.\,$n=0$でも成立する.

            格子点としての意味も確かめよう.\,$x=k$と固定したとき,
            \,$1\le y\le\frac{5k}{2}$を満たす整数$y$は$\lfloor\frac{5k}{2}\rfloor$個である.
            \,$k=1,\ldots,n$について足した値が$a_n$だから,【解答】の範囲の点を数えている.
            \bookterm{sequence}{数列}の値,\bookterm{floor-function}{床関数}の和,図形の中の個数が同じ量を表している.
    \end{enumerate}


  \medskip
  \noindent \textbf{【着眼点の振り返り】}\\
    ガウス記号の値を不等式で確定する方法,群や余りで値を整理する方法,
    添字を分割する方法,小数部分や個数として読む方法を見てきた.
    そのどれを使うかは,記号の位置と,式の中で果たす役割から考えよう.
    余りと桁,群数列と数え上げのように,一つの問題から別の見方へつながることもある.
